% !TEX root = ../main.tex
\clearpage
\begingroup
\setmonofont{Consolas}
\section*{附录B\quad 支撑材料文件清单}
支撑材料与论文分别提交。以下文件路径均相对于支撑材料根目录；根目录中的五个result工作簿为已填写的结果，templates目录中的同名文件为空白模板。

\begingroup
\small\setstretch{1.08}
\begin{enumerate}[label=\arabic*.,leftmargin=2.5em,itemsep=2pt]
\item \texttt{AI工具使用详情.pdf}
\item \texttt{code/\allowbreak{}energy\_\allowbreak{}paths.tex}
\item \texttt{code/\allowbreak{}explore\_\allowbreak{}data.py}
\item \texttt{code/\allowbreak{}export\_\allowbreak{}results.py}
\item \texttt{code/\allowbreak{}fig\_\allowbreak{}data.py}
\item \texttt{code/\allowbreak{}fig\_\allowbreak{}q1.py}
\item \texttt{code/\allowbreak{}fig\_\allowbreak{}q2\_\allowbreak{}sensitivity.py}
\item \texttt{code/\allowbreak{}fig\_\allowbreak{}q3\_\allowbreak{}days.py}
\item \texttt{code/\allowbreak{}fig\_\allowbreak{}q3\_\allowbreak{}quantile.py}
\item \texttt{code/\allowbreak{}framework.tex}
\item \texttt{code/\allowbreak{}prepare\_\allowbreak{}inputs.py}
\item \texttt{code/\allowbreak{}prepare\_\allowbreak{}source\_\allowbreak{}tables.py}
\item \texttt{code/\allowbreak{}q1.py}
\item \texttt{code/\allowbreak{}q2.py}
\item \texttt{code/\allowbreak{}q3.py}
\item \texttt{code/\allowbreak{}q3\_\allowbreak{}base.py}
\item \texttt{code/\allowbreak{}q4.py}
\item \texttt{code/\allowbreak{}q4\_\allowbreak{}tree.py}
\item \texttt{code/\allowbreak{}q4\_\allowbreak{}update.py}
\item \texttt{code/\allowbreak{}reproduce.py}
\item \texttt{code/\allowbreak{}source\_\allowbreak{}map.json}
\item \texttt{code/\allowbreak{}validation.py}
\item \texttt{data/\allowbreak{}actual\_\allowbreak{}power.csv.gz}
\item \texttt{data/\allowbreak{}actual\_\allowbreak{}price.csv.gz}
\item \texttt{data/\allowbreak{}day\_\allowbreak{}reference.csv.gz}
\item \texttt{data/\allowbreak{}inputs.npz}
\item \texttt{data/\allowbreak{}pv\_\allowbreak{}forecasts.csv.gz}
\item \texttt{data/\allowbreak{}q1\_\allowbreak{}config.json}
\item \texttt{data/\allowbreak{}raw\_\allowbreak{}hashes.json}
\item \texttt{evidence/\allowbreak{}export\_\allowbreak{}check.json}
\item \texttt{evidence/\allowbreak{}extra\_\allowbreak{}reproduction\_\allowbreak{}check.json}
\item \texttt{evidence/\allowbreak{}Q1\_\allowbreak{}storage\_\allowbreak{}sensitivity.csv}
\item \texttt{evidence/\allowbreak{}Q1\_\allowbreak{}strategy\_\allowbreak{}comparison.csv}
\item \texttt{evidence/\allowbreak{}Q2\_\allowbreak{}annual\_\allowbreak{}candidates.csv}
\item \texttt{evidence/\allowbreak{}Q2\_\allowbreak{}refinement.csv}
\item \texttt{evidence/\allowbreak{}Q2\_\allowbreak{}validation.csv}
\item \texttt{evidence/\allowbreak{}Q3\_\allowbreak{}final\_\allowbreak{}comparison.csv}
\item \texttt{evidence/\allowbreak{}Q4\_\allowbreak{}final\_\allowbreak{}comparison.csv}
\item \texttt{evidence/\allowbreak{}Q4\_\allowbreak{}validation.csv}
\item \texttt{evidence/\allowbreak{}reproduction\_\allowbreak{}check.json}
\item \texttt{evidence/\allowbreak{}workbook\_\allowbreak{}readback.json}
\item \texttt{evidence/\allowbreak{}参考文献来源.tex}
\item \texttt{figure\_\allowbreak{}inputs/\allowbreak{}derived\_\allowbreak{}values.json}
\item \texttt{figure\_\allowbreak{}inputs/\allowbreak{}eda/\allowbreak{}forecast\_\allowbreak{}error\_\allowbreak{}by\_\allowbreak{}lead.csv}
\item \texttt{figure\_\allowbreak{}inputs/\allowbreak{}eda/\allowbreak{}monthly\_\allowbreak{}summary.csv}
\item \texttt{figure\_\allowbreak{}inputs/\allowbreak{}eda/\allowbreak{}weekday\_\allowbreak{}summary.csv}
\item \texttt{figure\_\allowbreak{}inputs/\allowbreak{}q1/\allowbreak{}outputs/\allowbreak{}schedules/\allowbreak{}C03-p07.csv}
\item \texttt{figure\_\allowbreak{}inputs/\allowbreak{}q1/\allowbreak{}outputs/\allowbreak{}schedules/\allowbreak{}C04.csv}
\item \texttt{figure\_\allowbreak{}inputs/\allowbreak{}q1/\allowbreak{}outputs/\allowbreak{}schedules/\allowbreak{}C10-p05.csv}
\item \texttt{figure\_\allowbreak{}inputs/\allowbreak{}q1/\allowbreak{}outputs/\allowbreak{}summary.json}
\item \texttt{figure\_\allowbreak{}inputs/\allowbreak{}q2/\allowbreak{}refinement/\allowbreak{}comparison.csv}
\item \texttt{figure\_\allowbreak{}inputs/\allowbreak{}q2/\allowbreak{}supplement/\allowbreak{}comparison.csv}
\item \texttt{figure\_\allowbreak{}inputs/\allowbreak{}q2/\allowbreak{}validation/\allowbreak{}comparison.csv}
\item \texttt{figure\_\allowbreak{}inputs/\allowbreak{}q3\_\allowbreak{}daily.csv}
\item \texttt{figures/\allowbreak{}data\_\allowbreak{}compact.pdf}
\item \texttt{figures/\allowbreak{}forecast\_\allowbreak{}compact.pdf}
\item \texttt{figures/\allowbreak{}model.pdf}
\item \texttt{figures/\allowbreak{}q1\_\allowbreak{}1.png}
\item \texttt{figures/\allowbreak{}q1\_\allowbreak{}daily\_\allowbreak{}dispatch.pdf}
\item \texttt{figures/\allowbreak{}q3\_\allowbreak{}final\_\allowbreak{}days.pdf}
\item \texttt{figures/\allowbreak{}q3\_\allowbreak{}final\_\allowbreak{}quantile.pdf}
\item \texttt{figures/\allowbreak{}sensitivity.png}
\item \texttt{figures/\allowbreak{}weiwang.png}
\item \texttt{figures/\allowbreak{}图源说明.md}
\item \texttt{requirements.txt}
\item \texttt{result1.xlsx}
\item \texttt{result2.xlsx}
\item \texttt{result3.xlsx}
\item \texttt{result4-2.xlsx}
\item \texttt{result4-3.xlsx}
\item \texttt{results/\allowbreak{}daily\_\allowbreak{}costs.csv}
\item \texttt{results/\allowbreak{}Q1.csv}
\item \texttt{results/\allowbreak{}Q2.npz}
\item \texttt{results/\allowbreak{}Q3.npz}
\item \texttt{results/\allowbreak{}Q3\_\allowbreak{}adjustments.csv.gz}
\item \texttt{results/\allowbreak{}Q4-2.npz}
\item \texttt{results/\allowbreak{}Q4-3.npz}
\item \texttt{results/\allowbreak{}Q4-3\_\allowbreak{}adjustments.csv.gz}
\item \texttt{SHA256SUMS.txt}
\item \texttt{templates/\allowbreak{}result1.xlsx}
\item \texttt{templates/\allowbreak{}result2.xlsx}
\item \texttt{templates/\allowbreak{}result3.xlsx}
\item \texttt{templates/\allowbreak{}result4-2.xlsx}
\item \texttt{templates/\allowbreak{}result4-3.xlsx}
\item \texttt{使用说明.md}
\item \texttt{文件清单.md}
\end{enumerate}
\endgroup

\bigskip
\section*{附录C\quad 程序运行与结果文件说明}
\subsection*{C.1\quad 运行环境与入口}
计算程序使用Python 3.13，所需软件包及版本见支撑材料中的requirements.txt。程序使用相对路径读取文件，在支撑材料根目录依次执行以下命令：
\begin{Verbatim}[fontsize=\small,frame=single,xleftmargin=8pt,xrightmargin=8pt]
python -m pip install -r requirements.txt
python code/prepare_inputs.py
python code/reproduce.py
python code/export_results.py
\end{Verbatim}
第一条命令安装依赖；第二条根据随包数据重建计算输入并逐值核对；第三条重新计算问题一、衔接1月份储电量，再分别计算后续四条年度方案。复算输出写入reproduced目录，并与results目录的保存结果比较。最后一条命令依据results中的完整精度记录填写五份结果模板，核对费用和储能约束。

如需从赛题原始Excel重新整理数据，可将附件按原目录结构置于\texttt{data/\allowbreak{}raw/\allowbreak{}C题/\allowbreak{}附件/\allowbreak{}}，将题面置于\texttt{data/\allowbreak{}raw/\allowbreak{}C题/\allowbreak{}C题.pdf}，再运行\texttt{code/\allowbreak{}prepare\_\allowbreak{}source\_\allowbreak{}tables.py}。原始附件由赛题提供，不在支撑材料中重复保存。

\subsection*{C.2\quad 结果与图表的对应关系}
result1.xlsx给出问题一完整一天的安排；result2.xlsx、result3.xlsx、result4-2.xlsx与result4-3.xlsx分别给出对应方案2025年2月1日至12月31日的结果。工作簿沿用赛题模板样式，单元格保留完整精度，费用以未舍入值计算。results目录保存逐时段记录和逐日费用，问题三及问题四日内分支的历次调整另存于相应的压缩CSV文件。

\texttt{code/\allowbreak{}explore\_\allowbreak{}data.py}生成描述统计，\texttt{code/\allowbreak{}validation.py}提供已有对照检验的运行入口。五个以fig开头的Python程序根据figure\_inputs中的数据生成图表；两份TikZ源文件可使用XeLaTeX编译。具体命令见支撑材料的使用说明。本文结果表由Python写入Excel，不依赖Excel宏或SPSS交互命令。

\clearpage
\section*{附录D\quad 完整源程序}
以下按数据整理、模型计算、检验与结果导出、图表绘制的顺序列出支撑材料中的全部程序。程序内容直接读取同名源文件，不以片段或省略号代替。
\begingroup
\setstretch{1.0}

\par\ifdim\dimexpr\pagegoal-\pagetotal\relax<8\baselineskip\newpage\fi
\subsection*{D.1\quad 原始附件读取与整理}
\noindent 文件：\texttt{code/\allowbreak{}prepare\_\allowbreak{}source\_\allowbreak{}tables.py}\par
\smallskip
\VerbatimInput[breaklines=true,breakanywhere=true,fontsize=\footnotesize,baselinestretch=1,numbers=left,numbersep=4pt,xleftmargin=18pt,xrightmargin=6pt,frame=single,framesep=3pt]{../../../支撑材料/code/prepare_source_tables.py}

\par\ifdim\dimexpr\pagegoal-\pagetotal\relax<8\baselineskip\newpage\fi
\subsection*{D.2\quad 计算输入重建}
\noindent 文件：\texttt{code/\allowbreak{}prepare\_\allowbreak{}inputs.py}\par
\smallskip
\VerbatimInput[breaklines=true,breakanywhere=true,fontsize=\footnotesize,baselinestretch=1,numbers=left,numbersep=4pt,xleftmargin=18pt,xrightmargin=6pt,frame=single,framesep=3pt]{../../../支撑材料/code/prepare_inputs.py}

\par\ifdim\dimexpr\pagegoal-\pagetotal\relax<8\baselineskip\newpage\fi
\subsection*{D.3\quad 数据探索性分析}
\noindent 文件：\texttt{code/\allowbreak{}explore\_\allowbreak{}data.py}\par
\smallskip
\VerbatimInput[breaklines=true,breakanywhere=true,fontsize=\footnotesize,baselinestretch=1,numbers=left,numbersep=4pt,xleftmargin=18pt,xrightmargin=6pt,frame=single,framesep=3pt]{../../../支撑材料/code/explore_data.py}

\par\ifdim\dimexpr\pagegoal-\pagetotal\relax<8\baselineskip\newpage\fi
\subsection*{D.4\quad 问题一线性规划}
\noindent 文件：\texttt{code/\allowbreak{}q1.py}\par
\smallskip
\VerbatimInput[breaklines=true,breakanywhere=true,fontsize=\footnotesize,baselinestretch=1,numbers=left,numbersep=4pt,xleftmargin=18pt,xrightmargin=6pt,frame=single,framesep=3pt]{../../../支撑材料/code/q1.py}

\par\ifdim\dimexpr\pagegoal-\pagetotal\relax<8\baselineskip\newpage\fi
\subsection*{D.5\quad 问题二预测与跨日调度}
\noindent 文件：\texttt{code/\allowbreak{}q2.py}\par
\smallskip
\VerbatimInput[breaklines=true,breakanywhere=true,fontsize=\footnotesize,baselinestretch=1,numbers=left,numbersep=4pt,xleftmargin=18pt,xrightmargin=6pt,frame=single,framesep=3pt]{../../../支撑材料/code/q2.py}

\par\ifdim\dimexpr\pagegoal-\pagetotal\relax<8\baselineskip\newpage\fi
\subsection*{D.6\quad 运行回放与共享计算}
\noindent 文件：\texttt{code/\allowbreak{}q3\_\allowbreak{}base.py}\par
\smallskip
\VerbatimInput[breaklines=true,breakanywhere=true,fontsize=\footnotesize,baselinestretch=1,numbers=left,numbersep=4pt,xleftmargin=18pt,xrightmargin=6pt,frame=single,framesep=3pt]{../../../支撑材料/code/q3_base.py}

\par\ifdim\dimexpr\pagegoal-\pagetotal\relax<8\baselineskip\newpage\fi
\subsection*{D.7\quad 问题三日内滚动调度}
\noindent 文件：\texttt{code/\allowbreak{}q3.py}\par
\smallskip
\VerbatimInput[breaklines=true,breakanywhere=true,fontsize=\footnotesize,baselinestretch=1,numbers=left,numbersep=4pt,xleftmargin=18pt,xrightmargin=6pt,frame=single,framesep=3pt]{../../../支撑材料/code/q3.py}

\par\ifdim\dimexpr\pagegoal-\pagetotal\relax<8\baselineskip\newpage\fi
\subsection*{D.8\quad 电价预测与调度基础}
\noindent 文件：\texttt{code/\allowbreak{}q4.py}\par
\smallskip
\VerbatimInput[breaklines=true,breakanywhere=true,fontsize=\footnotesize,baselinestretch=1,numbers=left,numbersep=4pt,xleftmargin=18pt,xrightmargin=6pt,frame=single,framesep=3pt]{../../../支撑材料/code/q4.py}

\par\ifdim\dimexpr\pagegoal-\pagetotal\relax<8\baselineskip\newpage\fi
\subsection*{D.9\quad 问题四零点分支}
\noindent 文件：\texttt{code/\allowbreak{}q4\_\allowbreak{}tree.py}\par
\smallskip
\VerbatimInput[breaklines=true,breakanywhere=true,fontsize=\footnotesize,baselinestretch=1,numbers=left,numbersep=4pt,xleftmargin=18pt,xrightmargin=6pt,frame=single,framesep=3pt]{../../../支撑材料/code/q4_tree.py}

\par\ifdim\dimexpr\pagegoal-\pagetotal\relax<8\baselineskip\newpage\fi
\subsection*{D.10\quad 问题四日内分支}
\noindent 文件：\texttt{code/\allowbreak{}q4\_\allowbreak{}update.py}\par
\smallskip
\VerbatimInput[breaklines=true,breakanywhere=true,fontsize=\footnotesize,baselinestretch=1,numbers=left,numbersep=4pt,xleftmargin=18pt,xrightmargin=6pt,frame=single,framesep=3pt]{../../../支撑材料/code/q4_update.py}

\par\ifdim\dimexpr\pagegoal-\pagetotal\relax<8\baselineskip\newpage\fi
\subsection*{D.11\quad 四问完整复算入口}
\noindent 文件：\texttt{code/\allowbreak{}reproduce.py}\par
\smallskip
\VerbatimInput[breaklines=true,breakanywhere=true,fontsize=\footnotesize,baselinestretch=1,numbers=left,numbersep=4pt,xleftmargin=18pt,xrightmargin=6pt,frame=single,framesep=3pt]{../../../支撑材料/code/reproduce.py}

\par\ifdim\dimexpr\pagegoal-\pagetotal\relax<8\baselineskip\newpage\fi
\subsection*{D.12\quad 已有参数与信息对照}
\noindent 文件：\texttt{code/\allowbreak{}validation.py}\par
\smallskip
\VerbatimInput[breaklines=true,breakanywhere=true,fontsize=\footnotesize,baselinestretch=1,numbers=left,numbersep=4pt,xleftmargin=18pt,xrightmargin=6pt,frame=single,framesep=3pt]{../../../支撑材料/code/validation.py}

\par\ifdim\dimexpr\pagegoal-\pagetotal\relax<8\baselineskip\newpage\fi
\subsection*{D.13\quad 结果表填写与约束核验}
\noindent 文件：\texttt{code/\allowbreak{}export\_\allowbreak{}results.py}\par
\smallskip
\VerbatimInput[breaklines=true,breakanywhere=true,fontsize=\footnotesize,baselinestretch=1,numbers=left,numbersep=4pt,xleftmargin=18pt,xrightmargin=6pt,frame=single,framesep=3pt]{../../../支撑材料/code/export_results.py}

\par\ifdim\dimexpr\pagegoal-\pagetotal\relax<8\baselineskip\newpage\fi
\subsection*{D.14\quad 数据分析图}
\noindent 文件：\texttt{code/\allowbreak{}fig\_\allowbreak{}data.py}\par
\smallskip
\VerbatimInput[breaklines=true,breakanywhere=true,fontsize=\footnotesize,baselinestretch=1,numbers=left,numbersep=4pt,xleftmargin=18pt,xrightmargin=6pt,frame=single,framesep=3pt]{../../../支撑材料/code/fig_data.py}

\par\ifdim\dimexpr\pagegoal-\pagetotal\relax<8\baselineskip\newpage\fi
\subsection*{D.15\quad 问题一结果图}
\noindent 文件：\texttt{code/\allowbreak{}fig\_\allowbreak{}q1.py}\par
\smallskip
\VerbatimInput[breaklines=true,breakanywhere=true,fontsize=\footnotesize,baselinestretch=1,numbers=left,numbersep=4pt,xleftmargin=18pt,xrightmargin=6pt,frame=single,framesep=3pt]{../../../支撑材料/code/fig_q1.py}

\par\ifdim\dimexpr\pagegoal-\pagetotal\relax<8\baselineskip\newpage\fi
\subsection*{D.16\quad 问题二灵敏度图}
\noindent 文件：\texttt{code/\allowbreak{}fig\_\allowbreak{}q2\_\allowbreak{}sensitivity.py}\par
\smallskip
\VerbatimInput[breaklines=true,breakanywhere=true,fontsize=\footnotesize,baselinestretch=1,numbers=left,numbersep=4pt,xleftmargin=18pt,xrightmargin=6pt,frame=single,framesep=3pt]{../../../支撑材料/code/fig_q2_sensitivity.py}

\par\ifdim\dimexpr\pagegoal-\pagetotal\relax<8\baselineskip\newpage\fi
\subsection*{D.17\quad 问题三指定日期图}
\noindent 文件：\texttt{code/\allowbreak{}fig\_\allowbreak{}q3\_\allowbreak{}days.py}\par
\smallskip
\VerbatimInput[breaklines=true,breakanywhere=true,fontsize=\footnotesize,baselinestretch=1,numbers=left,numbersep=4pt,xleftmargin=18pt,xrightmargin=6pt,frame=single,framesep=3pt]{../../../支撑材料/code/fig_q3_days.py}

\par\ifdim\dimexpr\pagegoal-\pagetotal\relax<8\baselineskip\newpage\fi
\subsection*{D.18\quad 问题三参数比较图}
\noindent 文件：\texttt{code/\allowbreak{}fig\_\allowbreak{}q3\_\allowbreak{}quantile.py}\par
\smallskip
\VerbatimInput[breaklines=true,breakanywhere=true,fontsize=\footnotesize,baselinestretch=1,numbers=left,numbersep=4pt,xleftmargin=18pt,xrightmargin=6pt,frame=single,framesep=3pt]{../../../支撑材料/code/fig_q3_quantile.py}

\par\ifdim\dimexpr\pagegoal-\pagetotal\relax<8\baselineskip\newpage\fi
\subsection*{D.19\quad 电量关系示意图}
\noindent 文件：\texttt{code/\allowbreak{}energy\_\allowbreak{}paths.tex}\par
\smallskip
\VerbatimInput[breaklines=true,breakanywhere=true,fontsize=\footnotesize,baselinestretch=1,numbers=left,numbersep=4pt,xleftmargin=18pt,xrightmargin=6pt,frame=single,framesep=3pt]{../../../支撑材料/code/energy_paths.tex}

\par\ifdim\dimexpr\pagegoal-\pagetotal\relax<8\baselineskip\newpage\fi
\subsection*{D.20\quad 四问总体框架图}
\noindent 文件：\texttt{code/\allowbreak{}framework.tex}\par
\smallskip
\VerbatimInput[breaklines=true,breakanywhere=true,fontsize=\footnotesize,baselinestretch=1,numbers=left,numbersep=4pt,xleftmargin=18pt,xrightmargin=6pt,frame=single,framesep=3pt]{../../../支撑材料/code/framework.tex}
\endgroup
\endgroup
